\begin{lemma}[formulaGtOne]
\uses{VerticalIntegral}
\label{formulaGtOne}
\lean{Perron.formulaGtOne}
% at /home/runner/work/trough/trough/library/PrimeNumberTheoremAnd/PerronFormula.lean:1119.0-1159.66
\leanok
For $x>1$ and $\sigma>0$, we have
  $$
  \frac1{2\pi i}
  \int_{(\sigma)}\frac{x^s}{s(s+1)}ds =1-1/x.
  $$
\end{lemma}
\begin{proof}
\uses{vertIntBoundLeft,residuePull1,tendsto_rpow_atTop_nhds_zero_of_norm_gt_one,residuePull2,limitOfConstantLeft,contourPull3}
\leanok
Let $f(s) = x^s/(s(s+1))$. Then $f$ is holomorphic on $\C \setminus {0,-1}$.

First pull the contour from $(\sigma)$ to $(-1/2)$, picking up a residue $1$ at $s=0$.

Next pull the contour from $(-1/2)$ to $(-3/2)$, picking up a residue $-1/x$ at
  $s=-1$.

Then pull the contour all the way to $(\sigma')$ with $\sigma'<-3/2$.

For $\sigma' < -3/2$, the integral is bounded by
  $x^{\sigma'}\int_\R\frac{1}{|(1+t ^ 2)(2+t ^ 2)|^{1/2}}dt$.

Therefore $\int_{(\sigma')}\to 0$ as $\sigma'\to\infty$.

So pulling contours gives $\int_{(-3/2)}=0$.
\end{proof}
